\textbf{结论名称：} 阿氏圆模型（初中版）：\(PA = k \cdot PB\ (k>0,\ k \neq 1)\) 时，动点 \(P\) 的轨迹是一个圆（阿波罗尼斯圆）。

\textbf{适用条件：} 平面上给定两定点 \(A,B\)，给定常数 \(k>0\) 且 \(k \neq 1\)，动点 \(P\) 满足 \(PA = k \cdot PB\)。

\textbf{变量说明：} \(A,B\) 为定点，其间距离 \(AB = d\)；\(k\) 为已知的比例系数；\(C\) 为线段 \(AB\) 的内分点，满足 \(\frac{CA}{CB}=k\)；\(D\) 为外分点，满足 \(\frac{DA}{DB}=k\)（\(D\) 在 \(AB\) 或 \(BA\) 的延长线上）。

\textbf{核心公式：} 轨迹圆以 \(CD\) 为直径，圆心在直线 \(AB\) 上。若记 \(AB = d\)，则圆的半径为
\[
\boxed{R = \frac{k}{|1-k^{2}|}\,d}
\]
圆心距 \(A\) 点的有向距离为 \(\displaystyle \frac{k^{2}}{k^{2}-1}d\)。

\textbf{使用场景：} 中考数学中求解形如“\(PA + k \cdot PB\)”或“\(PA + \lambda PB\)”的最值问题。通过识别阿氏圆，将带系数线段相似转化为另一线段，进而利用“两点之间线段最短”或“点与圆上点的距离”求最值。

\textbf{简单例子：} 已知 \(A(0,0),\; B(4,0)\)，动点 \(P\) 满足 \(PA = 2PB\)。求 \(P\) 的轨迹圆。
由 \(k=2,\ d=4\)，内分点 \(C\) 满足 \(\frac{AC}{CB}=2 \Rightarrow C\bigl(\frac{8}{3},0\bigr)\)；
外分点 \(D\) 满足 \(\frac{AD}{DB}=2 \Rightarrow D(8,0)\)。
轨迹圆以 \(CD\) 为直径，圆心为 \(\bigl(\frac{16}{3},0\bigr)\)，半径 \(R = \frac{8}{3}\)。
验证：取 \(P\bigl(\frac{16}{3},\frac{8}{3}\bigr)\)，由两点间距离公式得 \(PA = \frac{8\sqrt{5}}{3}\)，\(PB = \frac{4\sqrt{5}}{3}\)，满足 \(PA = 2PB\)。